\begin{corollary}[4.8.11]
\label{IV.2.4.8.11-fr}
Sean $k$ un cuerpo, $X$ un $k$-preesquema, $K$ una extensión de $k$,
$\overline Z$ un subpreesquema cerrado de $X_{(K)}$; entonces existe un
menor cuerpo de definición de $\overline Z$.
\end{corollary}

\begin{proof}
En efecto, si $\overline{\mathscr J}$ es el ideal cuasicoherente de
$\mathcal O_{X_{(K)}}$ que define $\overline Z$, se deduce de (4.8.8) que
esto equivale a la existencia de un menor cuerpo de definición de
$\overline{\mathscr J}$, cuya existencia resulta de (4.8.9).
\end{proof}

\begin{corollary}[4.8.12]
\label{IV.2.4.8.12-fr}
Sean $k$ un cuerpo, $X$, $Y$ dos $k$-preesquemas, suponiendo $Y$ separado,
$K$ una extensión de $k$, $\overline f:X_{(K)}\to Y_{(K)}$ un $K$-morfismo.
Existe entonces un menor cuerpo de definición de $\overline f$.
\end{corollary}

\begin{proof}
En efecto, en virtud de (4.8.8), la existencia de tal cuerpo equivale a la de
un menor cuerpo de definición del grafo $\overline Z$ de $\overline f$; pero,
como $\overline Z$ es un subpreesquema \emph{cerrado} de
$(X\times_kY)_{(K)}$ (\textbf I, 5.4.3), la existencia de un menor cuerpo de
definición de $\overline Z$ resulta de (4.8.11).
\end{proof}

\begin{proposition}[4.8.13]
\label{IV.2.4.8.13-fr}
Bajo las hipótesis de \textup{(4.8.11) (resp. (4.8.9), (4.8.10), (4.8.12))},
supongamos además que $X$ sea de tipo finito sobre $k$ \textup{(resp.} que $X$ sea de
tipo finito sobre $k$ y $\mathscr F$ coherente, \textup{resp.} que $X$ sea de tipo
finito sobre $k$ y $\mathscr F$ y $\mathscr G$ coherentes, \textup{resp.} que $X$ e $Y$
sean de tipo finito sobre $k$\textup{)}. Entonces, el menor cuerpo de definición de
$\overline Z$ \textup{(resp.} de $\overline{\mathscr H}$, de $\overline u$, de $\overline f$\textup{)} es
una extensión de tipo finito de $k$.
\end{proposition}

\begin{proof}
Es posible limitarse a tratar el caso de (4.8.9) y suponer $X$ afín, puesto que
$X$ es unión finita de abiertos afines de tipo finito sobre $k$. Con las
notaciones de (4.8.6), todo se reduce a demostrar que, si $A$ es una álgebra de
tipo finito sobre $k$ y $M$ un $A$-módulo de tipo finito, existe una
subextensión $K'$ de $K$ que satisface las condiciones \emph{a)}, \emph{b)} o \emph{c)} y que es de
tipo finito sobre $k$ (toda subextensión de una extensión de tipo finito es de
tipo finito). Ahora bien, sea $(m_i)_{1\leq i\leq n}$ un sistema de generadores
del $A$-módulo $M$; como $A$ es noetheriano, lo mismo ocurre con $A_{(K)}$, y
por tanto $\overline N$ es un $A_{(K)}$-módulo de tipo finito, que admite un
sistema finito de generadores $n_j$ ($1\leq j\leq r$) tales que
$n_j=\sum_i(m_i\otimes1)b_{ij}$,
donde $b_{ij}\in A_{(K)}$;
cada $b_{ij}$ se escribe además
$b_{ij}=\sum_h a_{ijh}\otimes c_{ijh}$,
donde $a_{ijh}\in A$ y $c_{ijh}\in K$;
es claro que la extensión $K'$ de $k$ generada por los $c_{ijh}$ responde a
la cuestión.
\end{proof}

